% Informe de la práctica 03. En Overleaf: copiar este archivo como main.tex.
% Subir a la carpeta images/ de Overleaf:
%   figuras (resultados/graficas/): evidencia.png, integracion.png, resolucion.png, filtros.png, concurrencia.png
%   tablas  (resultados/tablas/):   tabla_entorno.tex, tabla_validacion.tex, tabla_estado.tex, tabla_frio.tex,
%           tabla_integracion.tex, tabla_resolucion.tex, tabla_filtros.tex, tabla_concurrencia.tex
% Tablas y figuras se generan con: python experimentos/graficas.py resultados/<corrida>
% La bibliografía es biblist.bib. Lo marcado con \pendiente{} se completa antes de la entrega.
% Compilar en el repo: cd informe && tectonic -X compile informe.tex

\documentclass[12pt]{article}
\usepackage[spanish,es-noshorthands,es-nodecimaldot]{babel}
\usepackage{natbib}
\usepackage{url}
\usepackage[utf8]{inputenc}
\usepackage[T1]{fontenc}
\usepackage{lmodern}
\usepackage{amsmath}
\usepackage{graphicx}
\graphicspath{{images/}{../resultados/graficas/}}   % Overleaf: images/; repo: ../resultados/graficas/
\usepackage{parskip}
\usepackage{fancyhdr}
\usepackage{booktabs}
\usepackage{float}
\usepackage{xcolor}
\usepackage{listings}
\usepackage{caption}
\usepackage[hidelinks]{hyperref}
\usepackage[compact]{titlesec}

\usepackage{vmargin}
\setmarginsrb{3 cm}{2.5 cm}{2.5 cm}{2.5 cm} {1 cm}{1.5 cm}{1 cm}{1.5 cm}

\title{Servicio web de convolución de imágenes con PyCUDA y Docker}
\author{Michael Lata \\ Jordy Romero \\ Andres Morocho}
\date{\today}

\makeatletter
\let\thetitle\@title
\let\theauthor\@author
\let\thedate\@date
\makeatother

\pagestyle{fancy}
\fancyhf{}
\lhead{Servicio web de convolución con PyCUDA y Docker}
\cfoot{\thepage}

\newcommand{\pendiente}[1]{\textcolor{red}{\textbf{[PENDIENTE: #1]}}}
% Tablas generadas por experimentos/graficas.py: en Overleaf en images/; en el repo en ../resultados/tablas/.
\newcommand{\tabla}[1]{%
  \IfFileExists{images/#1}{\input{images/#1}}{%
  \IfFileExists{../resultados/tablas/#1}{\input{../resultados/tablas/#1}}{%
  \IfFileExists{#1}{\input{#1}}{\pendiente{falta #1}}}}}
% Figuras generadas por el mismo script; si aún no existen, se marca el pendiente en su lugar.
\newcommand{\figura}[2]{%
  \IfFileExists{images/#2}{\includegraphics[width=#1\linewidth]{#2}}{%
  \IfFileExists{../resultados/graficas/#2}{\includegraphics[width=#1\linewidth]{#2}}{%
  \pendiente{falta #2}}}}
\newcommand{\repo}{\url{https://github.com/Whiteherobot/convolucion-servicio-pycuda}}
\newcommand{\repoAnterior}{\url{https://github.com/Whiteherobot/convolucion-pycuda}}

\lstset{basicstyle=\ttfamily\small, columns=fullflexible, frame=single, breaklines=true,
        keywordstyle=\bfseries, literate={á}{{\'a}}1 {é}{{\'e}}1 {í}{{\'i}}1
        {ó}{{\'o}}1 {ú}{{\'u}}1 {ñ}{{\~n}}1 {→}{{$\to$}}1 {×}{{$\times$}}1}

\setcounter{tocdepth}{1}

\begin{document}

\renewcommand{\tablename}{Tabla}
\renewcommand{\figurename}{Figura }

%%%%%%%%%%%%%%%%%%%%%%%%%%%%%%%%%%%%%%%%%%%%%%%%%%%%%%%%%%%%%%%%%%%%%%%%%%%%%%%%%%%%%%%%%
%%%%%%%%%%%%%  INICIO PORTADA  %%%%%%%%%%%%%%%%%%%%%%%%%%%%%%%%%%%%%%%%%%%%%%%%%%%%%%%%%%
%%%%%%%%%%%%%%%%%%%%%%%%%%%%%%%%%%%%%%%%%%%%%%%%%%%%%%%%%%%%%%%%%%%%%%%%%%%%%%%%%%%%%%%%%
\begin{titlepage}
    \centering
    \IfFileExists{images/logo.png}{\includegraphics[width=0.5\linewidth]{logo.png}}{}\\[1.0 cm]
    \textsc{\LARGE Universidad Politécnica Salesiana}\\[2.0 cm]
    \textsc{\Large CP}\\[0.5 cm]
    \textsc{\large Computación Paralela}\\[0.5 cm]
    \rule{\linewidth}{0.2 mm} \\[0.4 cm]
    { \huge \bfseries \thetitle}\\
    \rule{\linewidth}{0.2 mm} \\[1.5 cm]

 \begin{flushleft} \large
    \emph{Autores:}\\
    \theauthor
 \end{flushleft}

  \begin{flushleft} \large
    \emph{Profesores:}\\
    Ing. Gabriel Leon \\
  \end{flushleft}

  \begin{flushleft} \normalsize
    Repositorio: \repo\\
    Commit medido: \texttt{c892339} (mismo código que \texttt{f988bfc}; solo agrega resultados)\\
    Commit de entrega: \texttt{\pendiente{hash del commit etiquetado entrega}}\\
    Práctica 02: \repoAnterior, etiqueta \texttt{entrega}, commit \texttt{1c04426}
  \end{flushleft}

  \vfill

  {\large \thedate}

\end{titlepage}

%%%%%%%%%%%%%%%%%%%%%%%%%%%%%%%%%%%%%%%%%%%%%%%%%%%%%%%%%%%%%%%%%%%%%%%%%%%%%%%%%%%%%%%%%
%%%%%%%%%%%%%    FIN  PORTADA  %%%%%%%%%%%%%%%%%%%%%%%%%%%%%%%%%%%%%%%%%%%%%%%%%%%%%%%%%%
%%%%%%%%%%%%%%%%%%%%%%%%%%%%%%%%%%%%%%%%%%%%%%%%%%%%%%%%%%%%%%%%%%%%%%%%%%%%%%%%%%%%%%%%%

\tableofcontents
\vspace{1em}

%%%%%%%%%%%%%%%%%%%%%%%%%%%%%%%%%%%%%%%%%%%%%%%%%%%%%%%%%%%%%%%%%%%%%%%%%%%%%%%%%%%%%%%%%
%%%%%%%%%%%%%    INICIO   TEXTO   %%%%%%%%%%%%%%%%%%%%%%%%%%%%%%%%%%%%%%%%%%%%%%%%%%%%%%%
%%%%%%%%%%%%%%%%%%%%%%%%%%%%%%%%%%%%%%%%%%%%%%%%%%%%%%%%%%%%%%%%%%%%%%%%%%%%%%%%%%%%%%%%%
\section{Introducción}

En la práctica 02 se aplicaron tres filtros de convolución (media, nitidez y laplaciano) a imágenes en
escala de grises de 2000, 4000 y 6000 píxeles de lado, con un kernel CUDA propio ejecutado desde
Python mediante PyCUDA \cite{klockner2012,pycudadoc}. Esta práctica convierte ese mismo cálculo en un
servicio web: el usuario sube una imagen desde el navegador, elige el filtro y recibe el resultado
procesado en la GPU. El servicio se empaqueta en una imagen Docker con acceso a la GPU mediante NVIDIA
Container Toolkit \cite{nvidiactk} y se despliega en un equipo con GPU NVIDIA. Después se valida que
el resultado del servicio coincide con el de PyCUDA directo y se mide de punta a punta: cuánto cuesta
la GPU, cuánto el servicio y cuánto la comunicación, y qué cambia con la resolución, la máscara y la
concurrencia.

\begin{center}
\texttt{Navegador $\to$ Flask (gunicorn) $\to$ cola GPU $\to$ PyCUDA $\to$ GPU $\to$ PNG $\to$ navegador}
\end{center}

\section{Relación con la práctica 02}

El código reutilizado proviene de \repoAnterior, etiqueta \texttt{entrega} (commit \texttt{1c04426});
se copió del commit \texttt{d5ab309}, que es el mismo código reagrupado en carpetas. Sin cambios se
conservan, en \texttt{app/modelo/}: el kernel \texttt{kernels/convolucion.cu} (un píxel por hilo,
distribución en X, ceros fuera de la imagen, \texttt{float32}), las máscaras de media
$3\times3$, $9\times9$ y $21\times21$, nitidez $3\times3$ y laplaciano $3\times3$ con los mismos
coeficientes \cite{gonzalez2018}, la clase \texttt{ConvolucionPyCUDA} (contexto, compilación,
memoria, transferencias y eventos CUDA), la lectura y escritura de imágenes, el recorte a 0--255 solo
al final y el cálculo de errores con su tolerancia. Las imágenes de prueba son las mismas y se
comprueban con su hash (\texttt{datos/SHA256SUMS}).

Lo nuevo es la aplicación web, la gestión de solicitudes, el contenedor, el despliegue y la
evaluación. No se repite la comparación C/C++, CUDA C, Python y PyCUDA de la práctica 02.

\section{Arquitectura del servicio}

El código sigue una separación modelo--vista--controlador. \texttt{app/modelo/} contiene el cálculo
(lo reutilizado) y la gestión de solicitudes (\texttt{solicitudes.py}); \texttt{app/controlador/} las
rutas HTTP; \texttt{app/vista/} la página y su JavaScript. El mismo \texttt{procesar()} de
\texttt{app/modelo/procesamiento.py} lo usan el servicio y la ejecución directa
(\texttt{python -m app.directo}), de modo que la lógica HTTP queda separada de la convolución.

\begin{table}[H]
\centering
\small
\caption{Contrato HTTP del servicio.}\label{tab:contrato}
\begin{tabular}{p{4.2cm}p{10.3cm}}
\toprule
Ruta & Qué hace \\
\midrule
\texttt{GET /} & Interfaz: subir imagen, elegir filtro y máscara, ver original y resultado, descargar. \\
\texttt{POST /api/procesar} & \texttt{multipart/form-data}: \texttt{imagen}, \texttt{filtro}
(\texttt{media}, \texttt{nitidez}, \texttt{laplaciano}), \texttt{mascara} (3, 9 o 21; solo media),
\texttt{hilos} (opcional, 256). Responde el identificador, las dimensiones, hilos y bloques, los
tiempos en ms y las URL del resultado. \\
\texttt{GET /api/resultado/<id>.png} & Resultado recortado a 0--255, para ver o descargar. \\
\texttt{GET /api/resultado/<id>.npy} & Resultado \texttt{float32} sin recortar, para comparar. \\
\texttt{GET /api/salud} & GPU, \texttt{contexto\_ms}, \texttt{compilacion\_ms}, cola y límites. \\
\bottomrule
\end{tabular}
\end{table}

Los errores se devuelven como \texttt{\{"error": "mensaje"\}} con 400 (entrada inválida), 413
(archivo o imagen demasiado grande), 503 (cola llena o GPU no disponible) o 500 (falla del
procesamiento). Los límites son: archivo de hasta 40\,MB, imagen de hasta $6000\times6000$, formato
real PNG, JPEG, PGM/PPM, BMP o TIFF (se comprueba con Pillow, no por la extensión) y cola de 4. Cada
solicitud recibe un identificador aleatorio (\texttt{uuid4}) que nombra sus archivos de resultado,
así que ninguna sobrescribe ni recibe la imagen de otra; los resultados se borran a los 10 minutos.
La interfaz desactiva el botón mientras procesa para evitar solicitudes duplicadas.

\section{Integración con PyCUDA y gestión de solicitudes}

En PyCUDA el contexto CUDA queda asociado al hilo que lo crea, mientras que gunicorn atiende cada
solicitud en un hilo distinto. Por eso el servicio corre en un solo proceso
(\texttt{gunicorn -w 1 --threads 8}) con un único hilo trabajador dueño del contexto: al arrancar
crea \texttt{ConvolucionPyCUDA} (inicializa el contexto y compila el kernel con \texttt{nvcc}) y
después atiende una cola acotada en orden de llegada (FIFO). Los hilos HTTP validan la entrada,
encolan el trabajo y esperan su resultado. Como máximo una solicitud está en la GPU y cuatro esperan;
la siguiente recibe 503 de inmediato, lo que hace el rechazo medible. La inicialización se registra
aparte (\texttt{contexto\_ms} y \texttt{compilacion\_ms} en \texttt{/api/salud}); la compilación se
hace sin caché para que su tiempo sea el costo real de \texttt{nvcc}.

En cada solicitud, el trabajador decodifica la imagen a \texttt{float32}, copia la imagen y la máscara
a la GPU, lanza el kernel, copia el resultado de vuelta y codifica el PNG. Las copias y el kernel se
miden con eventos CUDA, sincronizando antes de leerlos; la memoria del anfitrión es paginable, y la
reserva y liberación en la GPU quedan fuera de esos tiempos, como en la práctica 02. Si la GPU no está
disponible al arrancar, el servicio responde 503 con el motivo y nunca calcula en la CPU.

\section{Docker y despliegue}

\begin{table}[H]
\centering
\small
\caption{Decisiones del \texttt{Dockerfile} (detalle en \texttt{despliegue/README.md}).}\label{tab:docker}
\begin{tabular}{p{5.2cm}p{9.3cm}}
\toprule
Decisión & Motivo \\
\midrule
Base \texttt{nvidia/cuda:13.4.3-devel-ubuntu24.04} fijada con su digest & PyCUDA compila el kernel al
arrancar y necesita \texttt{nvcc}; las variantes \texttt{runtime} y \texttt{base} no lo traen. El
digest no se mueve aunque la etiqueta cambie. \\
\texttt{python3-dev} y \texttt{build-essential} & \texttt{pip} compila PyCUDA al instalarlo. \\
Entorno virtual en \texttt{/opt/venv} & Ubuntu 24.04 no permite \texttt{pip} fuera de uno. \\
Usuario \texttt{servicio} sin privilegios & El servicio no necesita ser root. \\
\texttt{gunicorn -w 1 --threads 8}, sin depurador ni recarga & Un proceso: un contexto CUDA. \\
\texttt{HEALTHCHECK} sobre \texttt{/api/salud} & Sano solo cuando la GPU ya está cargada. \\
\texttt{.dockerignore} & Excluye \texttt{.git}, \texttt{.venv}, imágenes de prueba, resultados,
pruebas e informe. \\
\bottomrule
\end{tabular}
\end{table}

El controlador NVIDIA y CUDA de la imagen son componentes distintos. El controlador vive en el
anfitrión (\texttt{nvidia-smi}) y el contenedor accede a él a través de NVIDIA Container Toolkit
\cite{nvidiactk}; el toolkit CUDA con \texttt{nvcc} va dentro de la imagen. Para que funcione, el
controlador debe soportar una versión de CUDA igual o mayor que la de la imagen \cite{cudacompat}. Por
esa razón la imagen se construye y se mide en el mismo equipo: un controlador que solo soporte hasta
CUDA 13.3 no arranca esta imagen. Los comandos exactos son:

\begin{lstlisting}[language=bash]
docker build -t convolucion-servicio:p03 .
docker run --rm --gpus all convolucion-servicio:p03 nvidia-smi       # el contenedor ve la GPU
docker run -d --name conv-p03 --gpus all --memory=4g -p 8000:8000 convolucion-servicio:p03
\end{lstlisting}

No se contó con créditos de nube con GPU, así que el servicio se desplegó en la laptop de Jordy Romero,
con acceso desde otro equipo de la misma red: se comprobó desde un teléfono abriendo la interfaz y
\texttt{/api/salud} en \texttt{http://<IP de la laptop>:8000}. La imagen se construyó sin caché desde
un clon limpio; el contenedor quedó en estado \texttt{healthy} y una solicitud de media $9\times9$
sobre la imagen de 2000 píxeles se procesó en la GPU. La Tabla~\ref{tab:entorno} resume el equipo y
las versiones, registrados automáticamente al empezar la corrida de mediciones.

\begin{table}[H]
\centering
\small
\caption{Equipo de despliegue y medición (de \texttt{entorno.json} de la corrida).}\label{tab:entorno}
\tabla{tabla_entorno.tex}
\end{table}

\section{Validación}

Las pruebas de \texttt{pruebas/} (pytest) hablan con el servicio por HTTP, así que se ejecutaron contra
el contenedor. Cubren la matriz pequeña de la práctica 02 (laplaciano sobre $5\times5$, con valores
calculados a mano y bordes con ceros), una imagen de $1001\times777$ con 128, 256 y 512 hilos (tamaño,
bloques y cobertura hasta el último píxel), la equivalencia entre PyCUDA directo y el \texttt{.npy}
del servicio para los tres filtros (error absoluto máximo y medio antes del recorte, bajo la
tolerancia de la práctica 02), entradas inválidas (archivo que no es imagen, formato no permitido,
archivo o imagen demasiado grandes, filtro, máscara e hilos inválidos: código, mensaje y servicio
vivo) y aislamiento (dos solicitudes simultáneas reciben cada una su resultado). Con 12 solicitudes a
la vez se obtuvieron 5 respuestas 200 y 7 respuestas 503, y el servicio siguió respondiendo.

En la revisión se encontraron tres errores del servicio: una imagen truncada respondía 500 en vez de
400, la limpieza de resultados vencidos fallaba si dos hilos borraban el mismo archivo y la espera
de una solicitud no tenía tiempo máximo. Se corrigieron antes de medir: la imagen que no se puede
decodificar responde 400, la limpieza ignora los archivos que otro hilo ya borró y una solicitud espera
como máximo 120\,s antes de responder 503. Las dos pruebas que documentaban los errores pasaron a ser
pruebas normales. En la corrida de mediciones, las 27 pruebas pasaron contra el contenedor medido (\texttt{27 passed in
8.38s}, archivo \texttt{pytest.txt}).

La corrida de mediciones repite la equivalencia sobre la imagen real de $2000\times2000$ y guarda los
errores (Tabla~\ref{tab:validacion}): para cada filtro, la salida de PyCUDA directo en el anfitrión se
compara con el \texttt{.npy} del servicio, antes del recorte, con la misma función y tolerancia
($10^{-3}$) de la práctica 02. La Figura~\ref{fig:evidencia} muestra los PNG que devolvió el servicio.

\begin{table}[H]
\centering
\small
\caption{Error absoluto entre PyCUDA directo y el servicio, antes del recorte ($2000\times2000$).}\label{tab:validacion}
\tabla{tabla_validacion.tex}
\end{table}

\begin{figure}[H]
\centering
\figura{1.0}{evidencia.png}
\caption{Imagen original de $2000\times2000$ y resultado del servicio para los tres filtros. Abajo, el
recuadro rojo ampliado. El laplaciano se ve casi negro porque sus valores quedan cerca de cero o son
negativos y el PNG los recorta a 0--255.}\label{fig:evidencia}
\end{figure}

\section{Metodología}\label{sec:metodologia}

Se varía un factor a la vez con 256 hilos por bloque fijos (Tabla~\ref{tab:experimentos}). Cada
configuración tiene una solicitud de calentamiento, que se guarda pero no entra en las estadísticas,
y 5 repeticiones; se reportan la mediana y el coeficiente de variación (CV). El arranque en frío se
mide aparte: se reinicia el contenedor tres veces y se toma la primera solicitud de cada arranque. Toda
la corrida la ejecuta \texttt{experimentos/correr.py}, que antes verifica el hash de las imágenes, que
el servicio tenga GPU y anota el entorno; los fallos quedan en el CSV con su estado y en
\texttt{errores.log}, sin estimar valores.

\begin{table}[H]
\centering
\small
\caption{Experimentos.}\label{tab:experimentos}
\begin{tabular}{p{2.6cm}p{5.6cm}p{6.2cm}}
\toprule
Experimento & Configuración & Se compara \\
\midrule
Integración & Media $9\times9$, $4000\times4000$ & PyCUDA directo en el anfitrión, directo en un
contenedor de la misma imagen y a través del servicio \\
Resolución & Media $9\times9$; 2000, 4000 y 6000 & Kernel, procesamiento del servidor y total del cliente \\
Filtros y máscara & $4000\times4000$; media 3, 9 y 21; nitidez; laplaciano & Tiempos del servicio \\
Concurrencia & Media $9\times9$, $4000\times4000$; 1 y 2 clientes simultáneos, 5 rondas & Latencia,
espera en cola y solicitudes exitosas por segundo \\
\bottomrule
\end{tabular}
\end{table}

Cada tiempo tiene un alcance distinto. \emph{Kernel}, \emph{CPU$\to$GPU} y \emph{GPU$\to$CPU} se miden
con eventos CUDA. \emph{Decodificación}, \emph{codificación PNG}, \emph{espera en cola} y
\emph{servidor} (desde que llega la solicitud hasta tener el PNG, incluida la escritura del resultado)
usan el reloj monotónico del servidor. \emph{Cliente} va desde el envío hasta tener el PNG descargado,
con el reloj del cliente; nunca se restan marcas de tiempo de relojes distintos, solo duraciones. El
cliente corre en el mismo equipo que el servidor (\texttt{localhost}), así que la red no entra en las
mediciones. En concurrencia, el rendimiento es el número de solicitudes exitosas dividido por la
duración total de la ronda.

\section{Resultados}

La Tabla~\ref{tab:estado} cuenta las solicitudes correctas y con error de cada experimento.

\begin{table}[H]
\centering
\small
\caption{Estado de las solicitudes de la corrida (incluye calentamientos).}\label{tab:estado}
\tabla{tabla_estado.tex}
\end{table}

\subsection{Arranque en frío}

\begin{table}[H]
\centering
\small
\caption{Arranque en frío en ms: tiempo hasta que el servicio ve la GPU tras \texttt{docker restart},
contexto CUDA, compilación del kernel y primera solicitud, frente a la mediana estable.}\label{tab:frio}
\tabla{tabla_frio.tex}
\end{table}

Tras cada reinicio, el servicio tardó entre 1092 y 1673\,ms en estar listo con la GPU. De ese tiempo,
crear el contexto CUDA tomó entre 111 y 118\,ms y compilar el kernel entre 204 y 246\,ms; el resto
corresponde a levantar el contenedor y el servidor. La primera solicitud de cada arranque tardó entre
1354 y 1369\,ms en el servidor, igual que la mediana estable de la misma configuración (1368\,ms en la
Tabla~\ref{tab:integracion}). Es decir, el arranque en frío no penaliza la primera solicitud: el
contexto y la compilación se hacen al iniciar el servicio, antes de que el healthcheck lo marque como
listo, y no al llegar la primera imagen.

\subsection{Integración y contenedor}

\begin{table}[H]
\centering
\small
\caption{Media $9\times9$ sobre $4000\times4000$: medianas en ms por modo de ejecución.}\label{tab:integracion}
\resizebox{\linewidth}{!}{\tabla{tabla_integracion.tex}}
\end{table}

\begin{figure}[H]
\centering
\figura{0.95}{integracion.png}
\caption{Desglose del tiempo por modo. Cada segmento es la mediana de su componente; al final, la
mediana del total del cliente (o de la GPU en los modos directos).}\label{fig:integracion}
\end{figure}

La parte de GPU es prácticamente la misma en los tres modos: unos 9.4--9.6\,ms de CPU$\to$GPU,
2.1--2.25\,ms de kernel y 12.7--13.3\,ms de GPU$\to$CPU, con un total de 24.4--25.7\,ms. Llevar el
cálculo al servicio no lo hace más lento en la GPU. Lo que cambia es todo lo que rodea al cálculo: a
través del servicio la solicitud tarda 1368\,ms en el servidor y 1377\,ms en el cliente, unas 53
veces el tiempo de GPU. La diferencia entre servidor y cliente (unos 9\,ms) es el HTTP local y la
descarga del PNG. Los CV altos de los modos directos (79.6\,\% y 43.8\,\%) no reflejan un
comportamiento general: se deben a una o dos repeticiones con una copia GPU$\to$CPU anómala (84.9\,ms
en el anfitrión; 45.1 y 31.4\,ms en el contenedor, frente a unos 13\,ms en el resto). Por eso se
reporta la mediana.

\subsection{Resolución}

\begin{table}[H]
\centering
\small
\caption{Servicio, media $9\times9$: medianas en ms por resolución.}\label{tab:resolucion}
\resizebox{\linewidth}{!}{\tabla{tabla_resolucion.tex}}
\end{table}

\begin{figure}[H]
\centering
\figura{0.95}{resolucion.png}
\caption{Desglose del tiempo del servicio por resolución.}\label{fig:resolucion}
\end{figure}

Al pasar de $2000\times2000$ a $4000\times4000$ y $6000\times6000$ los píxeles crecen 4 y 9 veces. El
kernel crece en proporción similar (0.605, 2.25 y 4.98\,ms; 3.7 y 8.2 veces) y lo mismo la parte de
GPU (5.38, 25.6 y 55.8\,ms), dominada por las copias. La codificación PNG también escala con los
píxeles (312, 1281 y 2855\,ms) y es, con diferencia, el componente más grande: entre el 91\,\% y el
94\,\% del tiempo del servidor, mientras que el kernel no llega al 0.2\,\%. El total del cliente pasa
de 337 a 1405 y a 3138\,ms. En $6000\times6000$ el CV del cliente sube a 21\,\% porque una repetición
tardó 4833\,ms con una codificación normal (2855\,ms); el tiempo extra quedó fuera de la GPU, de la
decodificación y del PNG, y no se identificó su causa.

\subsection{Filtros y tamaño de máscara}

\begin{table}[H]
\centering
\small
\caption{Servicio, $4000\times4000$: medianas en ms por filtro.}\label{tab:filtros}
\resizebox{\linewidth}{!}{\tabla{tabla_filtros.tex}}
\end{table}

\begin{figure}[H]
\centering
\figura{0.95}{filtros.png}
\caption{Desglose del tiempo del servicio por filtro y máscara.}\label{fig:filtros}
\end{figure}

El kernel depende del tamaño de la máscara y no de sus coeficientes: los tres filtros de $3\times3$
(media, nitidez y laplaciano) tardan 0.51--0.52\,ms. Con la media, pasar de $3\times3$ a $9\times9$ y
$21\times21$ multiplica las operaciones por píxel por 9 y 49, pero el kernel solo crece 4.4 y 23
veces (0.512, 2.25 y 11.7\,ms), lo que indica que con máscaras pequeñas el kernel está limitado por
el acceso a memoria más que por el cálculo. Aun con la máscara de $21\times21$, el kernel es el
1.1\,\% del tiempo del servidor. El tiempo total lo decide de nuevo la codificación PNG, que varía
según el contenido de la imagen resultante y no según el costo del filtro: va de 570\,ms (nitidez) a
2095\,ms (media $3\times3$). Por eso la media $21\times21$, con el kernel más caro, tiene un total
menor (1041\,ms) que la media $3\times3$ (2195\,ms).

\subsection{Concurrencia}

\begin{table}[H]
\centering
\small
\caption{Media $9\times9$, $4000\times4000$: latencia del cliente (mediana y rango intercuartílico),
espera en cola y servidor en ms, y solicitudes exitosas por segundo (mediana de las rondas).}\label{tab:concurrencia}
\tabla{tabla_concurrencia.tex}
\end{table}

\begin{figure}[H]
\centering
\figura{1.0}{concurrencia.png}
\caption{Izquierda: desglose de la latencia con 1 y 2 clientes. Derecha: rendimiento de cada
ronda.}\label{fig:concurrencia}
\end{figure}

Todas las solicitudes terminaron bien (5/5 y 10/10). Con 2 clientes, en cada ronda una solicitud se
atendió de inmediato (unos 1360\,ms en el servidor) y la otra esperó en la cola entre 1315 y
1337\,ms, prácticamente lo que tarda la primera, y terminó en unos 2700\,ms. La mediana de 2050\,ms
y la espera de 658\,ms de la tabla mezclan esos dos grupos, como muestra el rango intercuartílico
(1388--2724\,ms). El rendimiento no mejora: 0.74 solicitudes exitosas por segundo con 1 cliente y
0.73 con 2. Duplicar los clientes duplica la duración de la ronda.

\section{Análisis}

\subsection{¿Qué se reutilizó de la práctica 02 y qué cambió al convertirla en servicio?}

Se reutilizaron sin cambios el kernel, las máscaras, la clase de PyCUDA, el manejo de imágenes y la
validación numérica. Cambió quién crea y usa el contexto CUDA: en la práctica 02 lo hacía el programa
principal en una sola ejecución; ahora lo crea una vez el hilo trabajador del servicio y lo mantiene
vivo entre solicitudes, que llegan por HTTP a hilos distintos y se serializan en una cola. Se
agregaron la validación de entradas, los identificadores por solicitud, la codificación PNG, los
límites y los códigos de error.

\subsection{¿Se mantiene la equivalencia numérica?}

Sí. En la imagen de $2000\times2000$ el error absoluto máximo y medio entre PyCUDA directo y el
\texttt{.npy} del servicio fue 0 en los tres filtros (Tabla~\ref{tab:validacion}): el servicio usa el
mismo kernel, la misma GPU, los mismos datos en \texttt{float32} y la misma configuración de lanzamiento,
por lo que el resultado es idéntico. Las únicas diferencias aparecen al codificar el PNG, que redondea
a 8 bits (hasta 0.5 por píxel) y recorta a 0--255; en el laplaciano esto elimina los valores
negativos, como se ve en la Figura~\ref{fig:evidencia}. Por eso la equivalencia se comprueba con el
\texttt{.npy} y no con el PNG.

\subsection{¿Cómo se reparte el tiempo entre GPU, preparación, servicio y comunicación?}

En $4000\times4000$ con media $9\times9$, de los 1377\,ms que ve el cliente, la GPU usa unos
25\,ms: 9.4\,ms en subir la imagen, 2.25\,ms en el kernel y 12.7\,ms en bajar el resultado. La
preparación en CPU es mayor: 40.5\,ms en decodificar la imagen y 1281\,ms en codificar el PNG. El
resto del servidor (validación, escritura del resultado) suma unas decenas de milisegundos y la
comunicación HTTP local unos 9\,ms. Las copias pesan más que el kernel porque los datos viajan en
\texttt{float32}: 64\,MB de entrada y 64\,MB de salida en $4000\times4000$. El patrón se repite en las
tres resoluciones: la codificación PNG ocupa más del 90\,\% del tiempo del servidor, así que el
servicio está limitado por la CPU y no por la GPU.

\subsection{¿Qué diferencias hay entre la ejecución directa y el servicio? ¿Se puede atribuir algo a Docker?}

Entre PyCUDA directo en el anfitrión y en el contenedor, las medianas difieren poco: 25.7 frente a
24.4\,ms de GPU, y el kernel 2.24 frente a 2.09\,ms. Ambos usan el mismo código y la misma GPU, pero
el anfitrión tiene CUDA 13.5 y la imagen CUDA 13.4, de modo que el kernel lo compilan versiones
distintas de \texttt{nvcc}; esa diferencia de unos 0.15\,ms no se puede atribuir a Docker. Con 5
repeticiones y valores atípicos en las copias, tampoco hay evidencia de una sobrecarga del contenedor
en la parte de GPU. La diferencia grande está entre ejecución directa y servicio (unos 25\,ms frente a
1377\,ms) y viene del trabajo que agrega el servicio (decodificar, codificar el PNG, escribir el
resultado y HTTP), no de Docker.

\subsection{¿Cómo afectan la resolución, la máscara y la concurrencia a la latencia, el rendimiento y la memoria?}

\textbf{Resolución.} La latencia crece casi en proporción al número de píxeles (337, 1405 y
3138\,ms para 2000, 4000 y 6000) porque todos sus componentes escalan así, sobre todo la codificación
PNG. \textbf{Máscara.} El kernel crece con la máscara (de 0.512 a 11.7\,ms), pero su efecto en la
latencia es menor que el del contenido de la imagen resultante sobre el PNG; el filtro más caro en la
GPU no es el más lento para el cliente. \textbf{Concurrencia.} Con 2 clientes la latencia de la
solicitud que espera se duplica (unos 2700\,ms) y el rendimiento se mantiene en unas 0.73
solicitudes/s. \textbf{Memoria.} No se midió la memoria por solicitud; por el tamaño de los datos en
\texttt{float32}, la entrada y la salida en GPU ocupan unos 32, 128 y 288\,MB en las tres
resoluciones, poco frente a los 8192\,MiB de la GPU. Con el servicio en reposo, el proceso usaba
160\,MiB de la GPU (\texttt{nvidia-smi\_anfitrion.txt}).

\subsection{¿Cómo condicionan el contexto CUDA, los trabajadores del servidor y la política de cola lo observado?}

Con un único contexto y un único hilo trabajador, la GPU procesa una solicitud a la vez: dos clientes
simultáneos no se procesan en paralelo, sino que uno espera en la cola. Con 2 clientes, la solicitud que llega
segunda espera entre 1315 y 1337\,ms en la cola, lo mismo que tarda en atenderse la primera, y el
rendimiento queda en 0.73 solicitudes/s frente a 0.74 con un cliente. Los 8 hilos del servidor HTTP
reciben las solicitudes en paralelo, pero solo encolan el trabajo; la cola admite 4 solicitudes en
espera además de la que se procesa, y las siguientes reciben 503, como en la prueba de 12 solicitudes
simultáneas (5 respuestas 200 y 7 respuestas 503). Como el cuello de botella es la codificación PNG en
CPU y no la GPU, aumentar el rendimiento requeriría paralelizar esa codificación (por ejemplo, con
varios procesos) más que agregar contextos CUDA.

\subsection{¿Qué límites del equipo, la red o el procedimiento restringen las conclusiones?}

Ver la Sección~\ref{sec:limitaciones}.

\section{Limitaciones}\label{sec:limitaciones}

\begin{itemize}
\item El servidor es una laptop: la GPU se comparte con la pantalla y puede bajar su frecuencia con
      la temperatura. Se midió conectada a la corriente y con las aplicaciones cerradas.
\item El cliente corre en el mismo equipo; las mediciones no incluyen una red real. El acceso desde
      otro equipo solo se comprobó de forma funcional.
\item La concurrencia se evaluó con 1 y 2 clientes; con un solo trabajador de GPU, más clientes solo
      alargan la cola hasta el rechazo con 503.
\item El servicio no tiene autenticación ni límite de solicitudes por cliente; solo se expone durante
      la demostración.
\item Se analiza una sola corrida con 5 repeticiones por configuración. Hubo valores atípicos en
      las copias GPU$\to$CPU de los modos directos y una repetición lenta en $6000\times6000$, cuyas
      causas no se identificaron; las medianas los atenúan, pero los CV de esas filas son altos.
\item No se midió la memoria de GPU ni de CPU por solicitud; solo se estimó por el tamaño de los datos.
\item La codificación PNG domina el tiempo del servicio, de modo que las latencias reflejan sobre todo
      el rendimiento de la CPU y del compresor, no el de la GPU.
\item El registro del servidor muestra cinco avisos de Gunicorn
      (\texttt{Control server error: Permission denied: '/home/servicio'}) porque el usuario del
      contenedor no tiene directorio personal. No afectaron a ninguna solicitud.
\end{itemize}

\section{Conclusiones}

\begin{enumerate}
\item El cálculo de la práctica 02 se convirtió en servicio sin perder equivalencia: el error frente a
      PyCUDA directo fue 0 en los tres filtros, y las 27 pruebas automáticas pasaron contra el
      contenedor medido.
\item Integrar el cálculo en el servicio y en Docker no cambia el tiempo de GPU (24.4--25.7\,ms en los
      tres modos). Las pequeñas diferencias entre anfitrión y contenedor no se pueden atribuir a
      Docker, porque cambian también la versión de CUDA.
\item El cuello de botella del servicio es la codificación PNG en CPU: ocupa más del 90\,\% del tiempo
      del servidor (1281 de 1368\,ms en $4000\times4000$), mientras que el kernel tarda 2.25\,ms. La
      latencia escala con los píxeles (337 a 3138\,ms) y depende más del contenido de la imagen
      resultante que del tamaño de la máscara.
\item Con un único contexto CUDA y un único trabajador, las solicitudes se atienden en serie: con 2
      clientes una espera unos 1.3\,s en la cola y el rendimiento se mantiene en 0.73 solicitudes/s.
      Para atender más clientes conviene paralelizar la codificación antes que la GPU.
\end{enumerate}

\section{Contribuciones}

\begin{table}[H]
\centering
\small
\begin{tabular}{p{3cm}p{11.2cm}}
\toprule
Integrante & Aporte \\
\midrule
Michael Lata & Aplicación web: servicio HTTP con Flask, contrato de la API, validación de entradas,
cola de la GPU con un único contexto CUDA, interfaz y ejecución directa (\texttt{app.directo}). \\
Jordy Romero & \texttt{Dockerfile} y \texttt{.dockerignore}, NVIDIA Container Toolkit, despliegue en su
laptop con acceso desde la red, construcción desde un clon limpio, pruebas con pytest y revisión del
servicio. Ejecución de las mediciones. \\
Andres Morocho & Protocolo de medición, cliente HTTP y scripts de los experimentos con verificación
del entorno, CSV, tablas y gráficas, informe y presentación. \\
\bottomrule
\end{tabular}
\end{table}

\section{Reproducibilidad}

El código está en \repo. Las mediciones se tomaron con el commit indicado en la portada, registrado
en \texttt{resultados/<corrida>/entorno.json} junto con el identificador de la imagen Docker. Las
imágenes de prueba se generan como indica \texttt{datos/README.md}. En un equipo con GPU NVIDIA,
Docker y NVIDIA Container Toolkit:

\begin{lstlisting}[language=bash]
git clone https://github.com/Whiteherobot/convolucion-servicio-pycuda && cd convolucion-servicio-pycuda
docker build -t convolucion-servicio:p03 .
docker run -d --name conv-p03 --gpus all --memory=4g -p 8000:8000 convolucion-servicio:p03
python3 -m venv .venv && .venv/bin/pip install -r requirements.txt
.venv/bin/python experimentos/correr.py        # resultados/<corrida>/
.venv/bin/python experimentos/graficas.py      # tablas y gráficas
\end{lstlisting}

\section*{Firmas}
\addcontentsline{toc}{section}{Firmas}
\vspace{1.5cm}
\noindent
\begin{tabular}{ccc}
\rule{4.3cm}{0.4pt} & \rule{4.3cm}{0.4pt} & \rule{4.3cm}{0.4pt} \\
Michael Lata & Jordy Romero & Andres Morocho \\
\end{tabular}

%%%%%%%%%%%%%%%%%%%%%%%%%%%%%%%%%%%%%%%%%%%%%%%%%%%%%%%%%%%%%%%%%%%%%%%%%%%%%%%%%%%%%%%%%
%%%%%%%%%%%%%     FIN   TEXTO     %%%%%%%%%%%%%%%%%%%%%%%%%%%%%%%%%%%%%%%%%%%%%%%%%%%%%%%
%%%%%%%%%%%%%%%%%%%%%%%%%%%%%%%%%%%%%%%%%%%%%%%%%%%%%%%%%%%%%%%%%%%%%%%%%%%%%%%%%%%%%%%%%

\bibliographystyle{plain}
\bibliography{biblist}

\end{document}
